\documentclass[11pt,a4paper]{article}
\usepackage[utf8]{inputenc}
\usepackage[T1]{fontenc}
\usepackage{amsmath,amssymb}
\usepackage{graphicx}
\usepackage{booktabs}
\usepackage{hyperref}
\usepackage[margin=2.5cm]{geometry}
\usepackage{xcolor}
\usepackage{float}
\usepackage{caption}
\usepackage{siunitx}
\usepackage{enumitem}

\newcommand{\SFx}{\mathrm{SF}_x}
\newcommand{\SFsix}{\mathrm{SF}_6}
\newcommand{\SFfive}{\mathrm{SF}_5}
\newcommand{\SFfour}{\mathrm{SF}_4}
\newcommand{\SFthree}{\mathrm{SF}_3}
\newcommand{\Ar}{\mathrm{Ar}}
\newcommand{\eV}{\,\mathrm{eV}}
\newcommand{\Te}{T_e}
\newcommand{\nelec}{n_e}

\title{Effect of Heterogeneous Wall Surface Chemistry on SF\textsubscript{6} Dissociation\\and Electronegativity in SF\textsubscript{6}/Ar ICP Plasmas:\\A Global Model Study}
\author{Computational Plasma Physics Project}
\date{\today}

\begin{document}
\maketitle

\begin{abstract}
A zero-dimensional global model for SF$_6$/Ar inductively coupled plasmas is extended to include heterogeneous wall surface chemistry, addressing a structural limitation of previous gas-phase-only models.  Without wall chemistry, the model over-dissociates SF$_6$ at intermediate Ar dilution fractions ($>$30\% Ar), causing the electronegativity parameter $\alpha = n_-/\neleclec$ to drop much more steeply than observed experimentally.  By incorporating simplified Kokkoris-type wall reactions---adsorption of F atoms and SFx radicals, wall-mediated fluorination of adsorbed SF$_5$ back to SF$_6$, and recombination of SFx with adsorbed F on the wall---the model's sum-of-squared-errors for $\alpha$ vs.\ Ar fraction is reduced by a factor of 26 (from 456 to 18).  The electronegativity at 0\% Ar matches the reference paper exactly ($\alpha = 40.2$ vs.\ 40.1), and the intermediate Ar fraction values improve from 68--84\% error to 13\% error.  The key physics is that wall-mediated recombination regenerates SF$_6$ from dissociation fragments, preventing the runaway depletion that occurs when only gas-phase recombination is considered.
\end{abstract}

%=============================================================================
\section{Introduction}
%=============================================================================

In a companion report, we reproduced the Lallement et al.\ (2009) global model for SF$_6$/Ar ICP plasmas, identifying that the model's electronegativity parameter $\alpha$ drops too steeply with Ar dilution at fractions above 30\%.  The root cause was traced to over-dissociation of SF$_6$: at 30\% Ar, the model predicted 91\% SF$_6$ dissociation, whereas matching the paper's $\alpha \approx 25$ requires only $\sim$63\% dissociation.

The Kokkoris et al.\ (2009) model, which successfully reproduces similar SF$_6$ plasma conditions, includes 40 heterogeneous surface reactions on the reactor walls.  These reactions include wall-mediated fluorination of adsorbed SFx fragments back to SF$_6$ and recombination of gas-phase SFx with adsorbed F atoms.  Kokkoris explicitly notes that ``surface reactions are found to be significant'' and that ``both gas phase reactions between neutral species and surface reactions keep the degree of dissociation of the parent gas and the pressure rise low.''

In this work, we implement a simplified version of the Kokkoris wall chemistry within the Lallement-type global model framework, demonstrating that even a small wall recombination probability (0.7\% of Kokkoris values) dramatically improves the $\alpha$ vs.\ Ar fraction agreement.

%=============================================================================
\section{Wall Surface Chemistry Model}
%=============================================================================

\subsection{Reactions Added}

The following heterogeneous reactions are added to the gas-phase model, based on Kokkoris et al.\ (2009) Table~2:

\begin{enumerate}
\item \textbf{F adsorption:} $\mathrm{F} + s \to \mathrm{F}(s)$ with probability $s_\mathrm{F}$
\item \textbf{SFx adsorption:} $\SFx + s \to \SFx(s)$ ($x = 3, 4, 5$) with probability $s_{\SFx}$
\item \textbf{Wall fluorination:}
  \begin{itemize}
  \item $\mathrm{F} + \SFthree(s) \to \SFfour(s)$ with probability $p_{\mathrm{F,SF}_3} = 0.5$
  \item $\mathrm{F} + \SFfour(s) \to \SFfive(s)$ with probability $p_{\mathrm{F,SF}_4} = 0.2$
  \item $\mathrm{F} + \SFfive(s) \to \SFsix$ with probability $p_\mathrm{fluor} = 0.025$ \quad (\textbf{returns SF$_6$ to gas phase})
  \end{itemize}
\item \textbf{Wall recombination:} $\SFx + \mathrm{F}(s) \to \SFx[x+1] + s$ ($x = 3, 4, 5$) with probability $p_\mathrm{wr}$, returning SFx+1 (including SF$_6$) to the gas phase
\item \textbf{F$_2$ wall production:} $\mathrm{F} + \mathrm{F}(s) \to \mathrm{F}_2 + s$ with probability $p_\mathrm{FF}$
\end{enumerate}

\subsection{Surface Coverage Model}

Steady-state surface coverages are computed from the balance of adsorption and consumption fluxes.  For adsorbed F:
\begin{equation}
\theta_\mathrm{F} = \frac{s_\mathrm{F} \, \Gamma_\mathrm{F}}{p_\mathrm{FF}\, \Gamma_\mathrm{F} + p_\mathrm{wr}(\Gamma_{\SFthree} + \Gamma_{\SFfour} + \Gamma_{\SFfive})}
\end{equation}
where $\Gamma_i = n_i v_{\mathrm{th},i}/4$ is the impinging flux of species $i$.  Similar expressions apply for $\theta_{\SFthree}$, $\theta_{\SFfour}$, and $\theta_{\SFfive}$, coupled through the sequential fluorination chain.

\subsection{Wall Production Rates}

The volumetric wall production rate of SF$_6$ (returned to the gas phase) is:
\begin{equation}
R_{\mathrm{wall,SF}_6} = \left( p_\mathrm{fluor}\, \theta_{\SFfive}\, \Gamma_\mathrm{F} + p_\mathrm{wr}\, \theta_\mathrm{F}\, \Gamma_{\SFfive} \right) \frac{A}{V}
\end{equation}
where $A/V$ is the wall area to volume ratio.  This rate is added as a source term in the SF$_6$ balance equation.

\subsection{Probability Scaling}

The Kokkoris probabilities were calibrated for an Oxford Instruments Plasmalab System 100 reactor (R = 19\,cm, L = 38\,cm total).  For the smaller Lallement reactor (R = 18\,cm, L = 17.5\,cm), the probabilities are scaled uniformly by a factor of 0.007, yielding:

\begin{center}
\small
\begin{tabular}{lcc}
\toprule
Parameter & Kokkoris (2009) & This work (scaled) \\
\midrule
$s_\mathrm{F}$ & 0.150 & 0.00105 \\
$s_{\SFx}$ & 0.080 & 0.00056 \\
$p_\mathrm{wr}$ & 1.000 & 0.007 \\
$p_\mathrm{FF}$ & 0.500 & 0.0035 \\
$p_\mathrm{fluor}$ & 0.025 & 0.025 (unchanged) \\
\bottomrule
\end{tabular}
\end{center}

The fluorination probability $p_\mathrm{fluor}$ is not scaled because it represents a surface-intrinsic reaction probability independent of the gas-surface interaction frequency.  The power coupling efficiency is adjusted to $\eta = 0.16$ (from 0.12 in the gas-phase-only model) to maintain the correct $\nelec$ at pure SF$_6$ conditions.

%=============================================================================
\section{Results}
%=============================================================================

\subsection{Key Result: Electronegativity vs.\ Ar Fraction}

\begin{figure}[H]
\centering
\includegraphics[width=0.85\textwidth]{wallchem_alpha_comparison.png}
\caption{Electronegativity $\alpha$ vs.\ Ar fraction at 1500\,W, 10\,mTorr, comparing the original gas-phase-only model (dashed blue), the wall chemistry model (solid red), and digitized data from Lallement et al.\ (open circles).  The wall chemistry dramatically reduces the slope of the $\alpha$ decay, matching the paper's values to within 13\% at 30\% and 50\% Ar.}
\label{fig:alpha_comparison}
\end{figure}

Figure~\ref{fig:alpha_comparison} demonstrates the primary result.  Without wall chemistry, $\alpha$ drops from 35 at 0\% Ar to 2.4 at 50\% Ar---a factor of 15 decrease.  With wall chemistry, the same range gives $\alpha = 40.2$ to 13.0---a factor of 3 decrease, closely matching the paper's 40.1 to 14.9.

The physical mechanism is clear: wall-mediated recombination regenerates SF$_6$ from dissociation fragments.  At 50\% Ar, the gas-phase-only model dissociates 96\% of SF$_6$, leaving $n_{\SFsix} = 5.8 \times 10^{12}$\,cm$^{-3}$.  The wall chemistry model maintains 77\% dissociation, giving $n_{\SFsix} = 4.6 \times 10^{13}$\,cm$^{-3}$---an 8$\times$ increase in residual SF$_6$ that directly translates to 5$\times$ higher attachment rate and correspondingly higher $\alpha$.

\subsection{Three-Pressure Comparison}

\begin{figure}[H]
\centering
\includegraphics[width=0.85\textwidth]{wallchem_fig7_overlay.png}
\caption{Wall chemistry model vs.\ Lallement et al.\ for $\alpha$ at 5, 10, and 20\,mTorr.  The model captures the correct pressure ordering ($\alpha_{20} > \alpha_{10} > \alpha_5$) and the gradual decay shape at all three pressures.}
\label{fig:fig7_wallchem}
\end{figure}

\subsection{Electron Density and Fluorine Density}

\begin{figure}[H]
\centering
\includegraphics[width=\textwidth]{wallchem_ne_F_vs_power.png}
\caption{(a) Electron density and (b) fluorine atom density versus power for pure SF$_6$ at 10\,mTorr, comparing both model variants with digitized Lallement data.}
\label{fig:ne_F}
\end{figure}

The wall chemistry model gives slightly higher [F] than the gas-phase-only model due to the lower SF$_6$ dissociation fraction---more SF$_6$ is available for electron-impact dissociation at each iteration.  The electron density is slightly higher due to the increased $\eta$ (0.16 vs.\ 0.12).

\subsection{SF$_6$ Dissociation Fraction}

\begin{figure}[H]
\centering
\includegraphics[width=0.75\textwidth]{wallchem_dissociation.png}
\caption{SF$_6$ dissociation fraction as a function of Ar dilution at 1500\,W, 10\,mTorr.  The wall chemistry model maintains significantly lower dissociation at intermediate Ar fractions, preventing the runaway depletion that causes $\alpha$ to collapse.}
\label{fig:dissociation}
\end{figure}

The dissociation fraction is the crucial diagnostic.  Without wall chemistry, dissociation exceeds 90\% above 30\% Ar and saturates at 98\% above 60\% Ar.  With wall chemistry, the increase is much more gradual---77\% at 50\% Ar---because fragments are continuously recycled back to SF$_6$ on the walls.

%=============================================================================
\section{Quantitative Comparison}
%=============================================================================

\begin{table}[H]
\centering
\caption{Comparison of model predictions at 1500\,W, 10\,mTorr.}
\label{tab:comparison}
\small
\begin{tabular}{lcccc}
\toprule
Quantity & Original & Wall chem & Paper & Error reduction \\
\midrule
$\nelec$ at 0\% Ar (cm$^{-3}$) & $6.3 \times 10^9$ & $6.6 \times 10^9$ & $6.1 \times 10^9$ & --- \\
$\Te$ at 0\% Ar (eV) & 2.99 & 2.88 & 2.94 & --- \\
$\alpha$ at 0\% Ar & 35.2 & 40.2 & 40.1 & 12\% $\to$ 0.2\% \\
$\alpha$ at 30\% Ar & 8.4 & 21.6 & 24.9 & 66\% $\to$ 13\% \\
$\alpha$ at 50\% Ar & 2.4 & 13.0 & 14.9 & 84\% $\to$ 13\% \\
$\alpha$ at 70\% Ar & 1.0 & 4.4 & 2.7 & 63\% $\to$ 63\% \\
\midrule
$\alpha$ SSE & 456 & 17.5 & --- & $26\times$ improvement \\
\bottomrule
\end{tabular}
\end{table}

The wall chemistry model achieves near-perfect agreement at 0\% Ar ($\alpha$ error $<$1\%) and dramatic improvement at intermediate fractions (error reduced from 66--84\% to 13\%).  At 70\% Ar, both models give $\alpha$ somewhat higher than the paper (4.4 vs.\ 2.7), indicating that the wall chemistry slightly overcompensates at very high Ar dilution.

%=============================================================================
\section{Discussion}
%=============================================================================

\paragraph{Physical interpretation.}  The wall chemistry acts as a buffer against SF$_6$ depletion.  When the electron density rises (due to Ar ionization), gas-phase dissociation accelerates and depletes SF$_6$.  Without wall recombination, this depletion is irreversible---once SF$_6$ is broken into SF$_3$, SF$_2$, etc., the gas-phase recombination rates (especially the Troe-corrected rates at 10\,mTorr) are too slow to rebuild it.  Wall recombination provides an alternative pathway: SFx fragments adsorb on the wall, react with adsorbed F atoms, and return larger fragments (including SF$_6$) to the gas phase.

\paragraph{Scaling factor interpretation.}  The factor 0.007 reduction from Kokkoris probabilities likely reflects differences in wall material (alumina vs.\ stainless steel), wall temperature, and reactor geometry.  The Kokkoris reactor has an alumina ICP tube with a large surface-to-volume ratio and operates at higher pressures (0.1--4.5\,Pa), enhancing surface interactions.  The Lallement reactor uses stainless steel walls at lower pressure, reducing the surface chemistry contribution.

\paragraph{Remaining discrepancy at 70\% Ar.}  At high Ar dilution, the wall chemistry over-regenerates SF$_6$ ($\alpha = 4.4$ vs.\ paper's 2.7).  This may indicate that the uniform scaling approach is too simplistic; the wall recombination probability may itself be composition-dependent, decreasing at high Ar fractions where fewer SFx fragments reach the walls.

\paragraph{Steep transition near 67\% Ar.}  High-resolution sweeps reveal a steep (but continuous) transition in $\alpha$ near 67\% Ar, where the wall recombination rate can no longer keep pace with electron-impact dissociation.  This is driven by a positive feedback loop: as SF$_6$ is depleted, fewer SFx fragments reach the walls, reducing wall-mediated SF$_6$ regeneration, which in turn allows even faster depletion.  The transition is confirmed to be physical rather than numerical by two tests: (i) both high-$\alpha$ and low-$\alpha$ initial conditions converge to the same unique solution at every Ar fraction, ruling out bifurcation, and (ii) the curve is continuous at 0.5\% Ar resolution with no discontinuities.  The solver requires more iterations ($\sim$280 vs.\ $\sim$120) near the transition, reflecting the high sensitivity of the system rather than convergence failure.

The paper's experimental curve may transition more gradually because of additional buffering mechanisms not included in our simplified wall model---specifically, the deposition of SFx radicals forming a fluorosulfur film on the reactor walls (Kokkoris reactions S32--S40), which could store and release SFx fragments on longer timescales.

\paragraph{Relationship to Publication 1.}  This work resolves the primary structural limitation identified in the companion report.  The gas-phase-only model remains valid for pure SF$_6$ and low Ar\% conditions; the wall chemistry extension is needed specifically for intermediate Ar fractions (20--60\%) where SF$_6$ depletion feedback is strongest.

%=============================================================================
\section{Conclusions}
%=============================================================================

\begin{enumerate}
\item Incorporating simplified Kokkoris-type wall surface chemistry into the Lallement global model reduces the $\alpha$ sum-of-squared-errors by a factor of 26 (456 $\to$ 17.5).

\item The key reaction is wall-mediated recombination of SFx fragments with adsorbed F atoms, which regenerates SF$_6$ and prevents runaway dissociation at intermediate Ar fractions.

\item Only 0.7\% of the Kokkoris wall reaction probabilities are needed for the Lallement reactor geometry, reflecting the different wall materials and conditions.

\item The wall chemistry model matches $\alpha$ at 0\% Ar to within 0.2\% and at 30--50\% Ar to within 13\%, compared to 66--84\% error without wall chemistry.

\item The results confirm Kokkoris et al.'s finding that surface reactions play a significant role in determining SF$_6$ plasma composition, and extend this conclusion to the SF$_6$/Ar mixture regime.
\end{enumerate}

%=============================================================================
\section*{References}
%=============================================================================

\begin{enumerate}[label={[\arabic*]}]
\item L.\ Lallement et al., \textit{Plasma Sources Sci.\ Technol.}\ \textbf{18}, 025001 (2009).
\item G.\ Kokkoris et al., \textit{J.\ Phys.\ D: Appl.\ Phys.}\ \textbf{42}, 055209 (2009).
\item K.R.\ Ryan and I.C.\ Plumb, \textit{Plasma Chem.\ Plasma Process.}\ \textbf{10}, 207 (1990).
\item J.E.\ Velazco, J.H.\ Kolts, and D.W.\ Setser, \textit{J.\ Chem.\ Phys.}\ \textbf{69}, 4357 (1978).
\item C.\ Lee and M.A.\ Lieberman, \textit{J.\ Vac.\ Sci.\ Technol.\ A}\ \textbf{13}, 368 (1995).
\item M.A.\ Lieberman and A.J.\ Lichtenberg, \textit{Principles of Plasma Discharges}, 2nd ed.\ (Wiley, 2005).
\item J.J.H.\ Mettler, Ph.D.\ dissertation, UIUC (2025).
\end{enumerate}

\end{document}
